% ==============================================================================
% Section 7: Analytical Derivation of First & Second-Order Greeks
% Target: Page 8 (Exactly 1 Page)
% ==============================================================================

\section{Analytical Derivation of First- and Second-Order Greeks}
\label{sec:greeks_derivation}

\subsection{Risk Sensitivities in Continuous-Time Hedging}

In financial engineering and trading desk management, dynamic hedging requires decomposing contingent claim price changes into partial derivatives with respect to state variables and model parameters, collectively termed the \textbf{Option Greeks} \citep{hull2018options}. Under Merton's continuous dividend framework, these derivatives must incorporate the discount factor $e^{-qT}$.

\subsection{First-Order Analytical Derivatives}

Differentiating the Merton pricing formulas (Equations~\eqref{eq:merton_call} and~\eqref{eq:merton_put}) yields:

\begin{enumerate}[leftmargin=1.5em, itemsep=1pt, parsep=0pt, label=\textbf{\arabic*.}]
    \item \textbf{Delta ($\bm{\Delta}$):} Measures the rate of change of option value with respect to the underlying spot price:
    \begin{align}
        \Delta_C &= \pdiff{C}{S_0} = e^{-qT} N(d_1) \in (0, \, e^{-qT}) \label{eq:delta_call} \\
        \Delta_P &= \pdiff{P}{S_0} = -e^{-qT} N(-d_1) = e^{-qT} \left( N(d_1) - 1 \right) \in (-e^{-qT}, \, 0) \label{eq:delta_put}
    \end{align}
    \item \textbf{Gamma ($\bm{\Gamma}$):} Measures the curvature of option price with respect to spot price (rate of change of Delta):
    \begin{equation}
        \Gamma = \frac{\partial^2 C}{\partial S_0^2} = \frac{\partial^2 P}{\partial S_0^2} = \frac{e^{-qT} \phi(d_1)}{S_0 \sigma \sqrt{T}} > 0
        \label{eq:gamma}
    \end{equation}
    where $\phi(z) = \frac{1}{\sqrt{2\pi}} e^{-z^2/2}$ is the standard normal probability density function. Gamma is identical for both calls and puts and peaks strictly at-the-money.
    \item \textbf{Vega ($\bm{\nu}$):} Measures sensitivity to changes in annualized asset volatility $\sigma$:
    \begin{equation}
        \nu = \pdiff{C}{\sigma} = \pdiff{P}{\sigma} = S_0 e^{-qT} \phi(d_1) \sqrt{T} > 0
        \label{eq:vega}
    \end{equation}
    \item \textbf{Theta ($\bm{\Theta}$):} Measures time decay, or rate of loss in option value as time elapses ($\dd t = -\dd T$):
    \begin{align}
        \Theta_C &= -\pdiff{C}{T} = -\frac{S_0 \sigma e^{-qT} \phi(d_1)}{2\sqrt{T}} + q S_0 e^{-qT} N(d_1) - r K e^{-rT} N(d_2) \label{eq:theta_call} \\
        \Theta_P &= -\pdiff{P}{T} = -\frac{S_0 \sigma e^{-qT} \phi(d_1)}{2\sqrt{T}} - q S_0 e^{-qT} N(-d_1) + r K e^{-rT} N(-d_2) \label{eq:theta_put}
    \end{align}
    \item \textbf{Rho ($\bm{\rho}$):} Measures sensitivity to shifts in the risk-free interest rate $r$:
    \begin{equation}
        \rho_C = \pdiff{C}{r} = K T e^{-rT} N(d_2) > 0, \qquad \rho_P = \pdiff{P}{r} = -K T e^{-rT} N(-d_2) < 0
        \label{eq:rho}
    \end{equation}
\end{enumerate}

\subsection{Second-Order Cross-Greeks: Vanna and Volga}

Institutional volatility trading desks monitor higher-order cross-sensitivities:
\begin{itemize}[leftmargin=1.5em, itemsep=1pt, parsep=0pt]
    \item \textbf{Vanna ($\frac{\partial \Delta}{\partial \sigma} = \frac{\partial \nu}{\partial S_0}$):} Quantifies how Delta shifts with volatility: $\text{Vanna} = -e^{-qT} \phi(d_1) \frac{d_2}{\sigma}$.
    \item \textbf{Volga / Vomma ($\frac{\partial \nu}{\partial \sigma} = \frac{\partial^2 V}{\partial \sigma^2}$):} Measures convexity of option value with respect to volatility: $\text{Volga} = \nu \frac{d_1 d_2}{\sigma}$.
\end{itemize}

\clearpage
